\section{Research Experience (Selected)}

\begin{enumerate}

    \item \textbf{BLADE --- Bilevel Low-rank Augmented-Lagrangian Erasure for LLM Unlearning}: Existing LLM unlearning methods struggle with robustness: unbounded forget losses degrade model coherence, fixed-weight balancing cannot adapt as retain difficulty shifts mid-training, and methods that work on one benchmark falter under scaling or repeated application. We proposed BLADE, a constrained bilevel framework whose three mechanisms give smooth, predictable control over the optimization landscape: a clamped-entropy forget loss whose gradient is exactly zero once a token reaches sufficient uncertainty; an asymmetric augmented Lagrangian that permanently ratchets retain protection after any violation; and a bilevel structure confined to LoRA adapters that repairs retain damage before each forgetting step. BLADE dominates across three benchmark families, improving average composite scores over the strongest baselines by $6\%$ on TOFU, $9\%$ on MUSE Books, and $7\%$ on KnowUndo, and it remains stable under $4\times$ scaling and $4$ sequential unlearning steps on MUSE News where the best competing method collapses entirely.\\
    Supervisor: \textit{Prof. Dongwon Lee and Prof. Ahmad Mousavi} \hfill Status: Accepted in EMNLP 2026 [Main Conference].

    \item \textbf{Amazon Nova AI Challenge 2026 --- Team Lion-0xA}: Co-leading Penn State's team as one of 10 global finalists. Our role as the Red Team involves designing autonomous multi-agent systems for web application security testing and user simulators for interacting with coding agents. Project title: \textit{``Autonomous Co-Design and Multi-Agent Red Teaming for Advancing Trusted Agentic AI.''} Supported by a \$250{,}000 Amazon grant and over \$1.5M in AWS credits.\\
    Faculty Adviser: \textit{Prof. Dongwon Lee} \hfill Status: Ongoing (Expected completion: Dec. 2026).

    \item \textbf{5G Access Control Formal Analysis}: Presented CoreScan, a formal analysis framework for verifying access control in 5G core networks. Uncovered five new classes of exploitable privilege escalation vulnerabilities using assume-guarantee reasoning and compositional verification. Findings were responsibly disclosed to GSMA.\\
    Supervisor: \textit{Prof. Syed Rafiul Hussain} \hfill Status: Published in IEEE S\&P 2025.

    \item \textbf{CRISPR-DIPOFF}: Developed interpretable deep learning models (RNN and Transformer-based) for CRISPR Cas-9 off-target prediction, using evolutionary algorithms for hyperparameter tuning and integrated gradients for model interpretation. Outperformed state-of-the-art and extended established biological hypotheses on the sgRNA seed region.\\
    Supervisor: \textit{Prof. M. Sohel Rahman} \hfill Status: Published in \textit{Briefings in Bioinformatics, 2024}.

    \item \textbf{NIDS}: Proposed a semi-supervised, adversarially robust deep learning approach for network intrusion detection, combining K-Means clustering, Autoencoders, Decision Trees, and MLPs. Validated robustness against FGSM and GAN-based IDSGAN attacks.\\
    Supervisor: \textit{Prof. Md. Shohrab Hossain} \hfill Status: Published in \textit{NSysS, 2023}.

\end{enumerate}

\endinput
